% =============================================================================
% Abstrakt / Kurzfassung
% =============================================================================
\chapter*{Kurzfassung}
\addcontentsline{toc}{chapter}{Kurzfassung}

Die Auswertung fahrzeugspezifischer Daten ermöglicht es, bei kontinuierlichen Nutzungsmustern das Bewegungsverhalten vorherzusagen. Dadurch kann bidirektionales Laden effizienter eingesetzt werden, was sowohl die Fahrzeughalter als auch das Stromnetz entlastet. Ziel ist es, durch Machine Learning ein Modell zu entwerfen, das möglichst genau individuelles Verkehrsverhalten prognostizieren kann. Zu bewerkstelligen ist dieses Vorhaben durch eine Ablationsstudie; Es werden auf Basis einer Datengrundlage Modelle erstellt und untereinander verglichen, wodurch die Annahme entsteht, dass ein Modell besser abschneiden wird, als die anderen. Das stärkste Modell wird dann weiter entwickelt um die Ergebnisse zu optimieren und somit die Zuverlässigkeit und den Beitrag zum Stromnetz von elektrischen Fahrzeugen zu erhöhen.


\vspace{1cm}

\textbf{Schlüsselwörter:} Künstliche Intelligenz, Machine Learning, Prognose, Fahrzeugnutzung, Elektromobilität, Bidirektionales Laden, Random Forest, LightGBM
